\documentclass{article}
\usepackage{graphicx}
\usepackage[margin=1.5cm]{geometry}
\usepackage{amsmath}

\begin{document}
\twocolumn

\title{Chapter 1 Warm-Up: Python Lists and Tuples}
\author{Prof. Jordan C. Hanson}

\maketitle

\section{Python Lists}

\begin{enumerate}

\item A Python \textbf{list} is an ordered collection of objects.  Lists are written using square brackets, \texttt{[ ]}, with the elements separated by commas.

\begin{itemize}

\item Begin by creating a list containing three student names:

\begin{verbatim}
students = ["Alice","Bob","Carlos"]
\end{verbatim}

\item Print the entire list:

\begin{verbatim}
print(students)
\end{verbatim}

\item Individual elements of a list can be accessed using their position, or \textit{index}.  Python begins counting at zero.  Try:

\begin{verbatim}
print(students[0])
print(students[1])
print(students[2])
\end{verbatim}

\item Which student is stored at index \texttt{0}?  Which student is stored at index \texttt{2}?

\end{itemize}
\end{enumerate}


\section{Lists Inside Lists}

\begin{enumerate}

\item The elements of a list do not have to be simple strings or numbers.  A list can itself contain another list.

Suppose we want to represent a course together with its students.

\begin{itemize}

\item Create a list containing the name of a course and the list of students:

\begin{verbatim}
course = ["Python Programming",students]
\end{verbatim}

\item Print the complete object:

\begin{verbatim}
print(course)
\end{verbatim}

\item The first element of \texttt{course} is the course name:

\begin{verbatim}
print(course[0])
\end{verbatim}

\item The second element is itself a list containing the students:

\begin{verbatim}
print(course[1])
\end{verbatim}
\vspace{0.25cm}

\item Because \texttt{course[1]} is a list, we can select an element from that list as well.  Try:

\begin{verbatim}
print(course[1][0])
print(course[1][2])
\end{verbatim}

\end{itemize}
\end{enumerate}


\section{Python Tuples}

\begin{enumerate}

\item A Python \textbf{tuple} is another ordered collection of objects.  Tuples are similar to lists, but are written using parentheses, \texttt{( )}.

\begin{itemize}

\item Create two simple course lists:

\begin{verbatim}
course1 = ["Python Programming",["Alice", "Bob"]]
course2 = ["Calculus I",["Diana", "Ethan"]]
\end{verbatim}

\item Combine the two course lists into a tuple:

\begin{verbatim}
courses = (course1,course2)
\end{verbatim}

\item Print the complete tuple:

\begin{verbatim}
print(courses)
\end{verbatim}

\item Access the first and second course using:

\begin{verbatim}
print(courses[0])
print(courses[1])
\end{verbatim}
\vspace{0.25cm}

\item Print only the name of the second course:

\begin{verbatim}
print(courses[1][0])
\end{verbatim}
\vspace{0.25cm}

\item Print the name of the first student enrolled in the second course:

\begin{verbatim}
print(courses[1][1][0])
\end{verbatim}
\vspace{0.25cm}

\end{itemize}
\end{enumerate}


\section{Warm-Up Exercise}

Construct a Python object representing two laboratory sections.  Each laboratory section should contain a course name and a list of three students. Use the following information:

\begin{verbatim}
Physics Laboratory
  Maya
  Noah
  Olivia

Computer Laboratory
  Liam
  Sophia
  Ethan
\end{verbatim}

\end{document}